\section{Methodology: design}

\subsection{Scope and data sources}

The system answers questions about the quarterly results of US public companies. It uses three
public datasets from the SEC and one benchmark (\cref{tab:data}).

\begin{table}[ht]
\centering\small
\begin{tabularx}{\linewidth}{lXl}
\toprule
Source & What we use & Size \\
\midrule
\code{companyfacts.zip} & XBRL figures of every filer; we keep 10 metrics per 10-Q & 1.4 GB, 20,609 files \\
\code{submissions.zip} & List of filings per company, to find each 10-Q document & 1.5 GB \\
10-Q documents (EDGAR) & Text passages from the MD\&A section, for ``why'' questions & 140,226 filings \\
FinanceBench~\cite{financebench} & 150 questions with analyst answers, 84 PDF filings & 160 MB \\
\bottomrule
\end{tabularx}
\caption{Data sources.}
\label{tab:data}
\end{table}

XBRL (eXtensible Business Reporting Language) is the machine-readable version of the financial
statements that companies file together with each report. Each fact has a tag (for example
\code{Revenues}), a value, a unit and the start and end date of its period. This gives us exact
numbers without parsing tables. We keep ten metrics: revenue, cost of sales, gross profit, SG\&A,
R\&D, operating income, pretax income, net income, diluted EPS and operating cash flow. Reasons
(``why did sales fall?'') are not in XBRL, so they come from the text of the
\emph{Management's Discussion and Analysis} (MD\&A) section of each 10-Q.

\subsection{Answer workflow}

\Cref{fig:workflow} shows the five steps of a question. Only two of them call an LLM (shaded).

\begin{figure}[ht]
\centering
\begin{tikzpicture}[
  box/.style={draw,rounded corners=2pt,align=center,font=\small,minimum height=12mm,text width=25mm},
  llm/.style={box,fill=llmbox},
  path/.style={draw,rounded corners=2pt,align=center,font=\footnotesize,text width=40mm,minimum height=15mm},
  arr/.style={-{Stealth[length=2mm]},thick}]
  \node[llm] (parse) at (0,0)    {1. Parse\\\footnotesize question $\to$ JSON};
  \node[box] (check) at (3.6,0)  {2. Check\\\footnotesize company, metric, scope};
  \node[box] (find)  at (7.2,0)  {3. Find\\\footnotesize company and filing (SQL)};
  \node[box] (route) at (10.8,0) {4. Route\\\footnotesize by question type};
  \node[path] (num) at (0.9,-2.6)  {\textbf{Numbers}\\SQL templates over XBRL; change computed in SQL};
  \node[path] (why) at (5.4,-2.6)  {\textbf{Why}\\hybrid search in MD\&A; LLM writes claims that cite passages};
  \node[path] (doc) at (9.9,-2.6)  {\textbf{Document QA}\\search one filing; code checks inputs and evaluates the formula};
  \node[llm,text width=105mm] (write) at (5.4,-5.0) {5. Write\\\footnotesize LLM phrases the verified facts; code checks numbers and references};
  \draw[arr] (parse)--(check); \draw[arr] (check)--(find); \draw[arr] (find)--(route);
  \draw[thick] (route.south) -- ++(0,-0.45) coordinate (b);
  \draw[thick] (b) -- (b -| num.north);
  \draw[arr] (b -| num.north) -- (num.north);
  \draw[arr] (b -| why.north) -- (why.north);
  \draw[arr] (b -| doc.north) -- (doc.north);
  \draw[arr] (num.south) -- (num.south |- write.north);
  \draw[arr] (why.south) -- (write.north);
  \draw[arr] (doc.south) -- (doc.south |- write.north);
\end{tikzpicture}
\caption{The answer workflow. Shaded steps call the LLM; all other steps are code. The ``why''
path also calls the LLM to write its claims.}
\label{fig:workflow}
\end{figure}

\paragraph{Step 1: parse.} The LLM turns the question into a small JSON object: company (as written
in the question), metric, period (for example \code{Q2 2023}), question type (numeric, explanatory
or both), comparison (year over year, quarter over quarter or none) and search terms. The output must
match a fixed schema.

\paragraph{Step 2: check.} Code then checks the parse. The company must appear in the question
(``ask, don't guess''), the metric must be one of the ten, and a few scope rules always apply: Q4
and full-year numbers are only in the 10-K, and segment numbers are not in our figures. A question
without a ``why'' word is never sent to the explanation path. These checks catch the most common
mistakes of small models.

\paragraph{Step 3: find.} SQL looks up the company (by ticker, a known alias, or a normalised name)
and the filing that covers the period.

\paragraph{Step 4: route.} The question goes to one of three paths.
\begin{itemize}
  \item \textbf{Numbers.} One of the ten metrics for a quarter, or its change. A fixed SQL template
        reads the figure. For a year-over-year change it reads the prior-year column of the
        \emph{same} 10-Q, so restated values are used (for example Johnson \& Johnson after the
        Kenvue spin-off). The change in percent and in dollars is computed in SQL. Each figure cites
        the filing, statement, line label, XBRL tag and period.
  \item \textbf{Why.} Hybrid search (BM25 plus semantic, merged by reciprocal rank fusion) over the
        MD\&A passages of that one filing. If the best match is too weak (cosine similarity below
        0.45), the system declines. Otherwise the LLM writes short claims, each of which must cite a
        retrieved passage by its id; claims that cite an unknown passage are dropped.
  \item \textbf{Document QA.} Everything else: ratios, full years, other line items and text
        questions, as in FinanceBench. The system searches the filing of that period (top 10
        passages). For a calculation the LLM does not give the result. It names each input exactly
        as printed, with the passage it comes from, and a formula. Code checks that each input
        really is printed in that passage and then evaluates the formula (only arithmetic is
        allowed).
\end{itemize}

\paragraph{Step 5: write.} A second LLM call turns the verified facts into a short answer. It gets
only the question and the facts, each with a reference such as
\code{[F1] 3M, net sales, Q2 2023: \$8.325 billion}. Code rejects the answer if it contains a number
that is not in the facts or a reference that does not exist. In that case the system uses a fixed
template answer instead. An example answer:

\begin{lstlisting}
3M's net sales fell year over year in the second quarter of 2023. Net sales were
$8.325 billion in Q2 2023, compared with $8.702 billion in Q2 2022 [F1][F2]. The
computed change is a decrease of 4.3%, or $377 million.
\end{lstlisting}

The design rule is that the user never sees a number that code did not find or compute. A number
question takes two LLM calls (parse, write), a ``why'' question three (parse, claims, write) and a
document QA question two (parse, answer).

\subsection{Database design}

All data lives in one Postgres 16 database with the pgvector extension (\cref{tab:db}).

\begin{table}[ht]
\centering\small
\begin{tabularx}{\linewidth}{lrX}
\toprule
Table & Rows & Contents \\
\midrule
\code{companies}   & 9,695     & CIK, name, main ticker, normalised name for lookups \\
\code{documents}   & 140,520   & One row per filing: form, period, filing date, link, whether text is indexed \\
\code{figures}     & 3,022,763 & Metric, XBRL tag, line label, period, value, unit, derived flag \\
\code{metric\_tags}& 96,950    & Which XBRL tag each company uses for each metric \\
\code{chunks}      & 4,550,888 & Passages of about 180 words, with section, heading and a 384-dimension embedding \\
\bottomrule
\end{tabularx}
\caption{Database tables after the full load.}
\label{tab:db}
\end{table}

Two rules for the figures matter for correctness. A quarter is a fact whose period is about three
months long, labelled from its own end date, and the tag is chosen per filing from an ordered list
of candidates (companies use different tags for ``revenue''). Cash flow is only reported year to
date, so a quarter is computed as the year-to-date value minus the previous year-to-date value and
marked as derived. Every search is filtered to one filing first, so the 4.5 million passages do not
slow a question down.

\subsection{Cloud architecture}

We run everything on one AWS EC2 instance (c7i.8xlarge: 32 vCPUs, 64 GB memory, 250 GB disk,
Ubuntu, region us-east-1) with Docker Compose. The containers are Postgres with pgvector, the
FastAPI service, a Streamlit web page, the Spark jobs and, optionally, Ollama to serve a local LLM
on the CPU. S3 keeps a long-term copy of the raw downloads, the Parquet files and the passages.

We considered an EMR Spark cluster but chose one large machine. The slowest part of the pipeline is
downloading 140,000 documents, and the SEC allows at most 10 requests per second from one user, no
matter how many machines send them. The CPU-heavy part (splitting and embedding text) is spread
over the cores of one machine by Spark just as well as over several machines, without the cost of
starting a cluster. The instance is stopped when it is not used, and the heavy steps are resumable.

\subsection{Baselines}

To show what the workflow adds, we compare it with two simpler systems that use the same LLM.
\begin{itemize}
  \item \textbf{C: simple RAG.} The same passage index and hybrid search over the correct filing,
        top 10 passages, one LLM call that must cite passages. No parsing, routing, SQL, code
        arithmetic or checks.
  \item \textbf{N: closed book.} The LLM gets only the question and must answer from its own
        knowledge, with the same instruction to cite sources or say it does not know. This shows
        whether retrieval is needed at all.
\end{itemize}
We also run \textbf{A0}, our workflow with the template answer instead of the writer, to measure
what the writer LLM adds. A fourth baseline, giving the LLM the whole filing, was built but not run:
it mostly measures how much context a model can take and was the most expensive cell.
